\chapter{Considerações Finais e Perspectivas}
\label{chap:conclusoes}

\section{Síntese da Trajetória Biofísica e Metrológica}
\label{sec:sintese_biofisica}

A trajetória da avaliação da qualidade de imagem em tomografia computadorizada reflete a constante superação de abstrações matemáticas lineares simplificadas em direção à modelagem da complexidade biofísica:
\begin{enumerate}
  \item \textbf{A Fase Linear Analítica (1950--1990):} Partiu da teoria de detecção de sinais e do Observador Ideal Bayesiano, introduzindo filtros oculares (NPWE) para simular as limitações fisiológicas humanas em fundos uniformes;
  \item \textbf{A Modelagem Cortical (1990--2015):} Desenvolveu o Observador de Hotelling Canalizado (CHO), utilizando canais de frequência inspirados na arquitetura do córtex visual primário para superar o ruído estrutural de fundos anatômicos;
  \item \textbf{A Ruptura da Linearidade (2015--2026):} A introdução clínica de algoritmos de reconstrução por aprendizado profundo (DLR) quebrou as premissas de invariância espacial e estacionariedade, expondo os limites dos modelos analíticos tradicionais;
  \item \textbf{O Paradigma da Inteligência Artificial Perceptual e Otimização Multiobjetivo (2026+):} A emergência dos observadores computacionais por aprendizado profundo com mecanismos de auto-atenção (DLMO), calibrados diretamente contra a percepção de radiologistas e integrados a espaços de decisão tridimensionais $(D, T, -W)$.
\end{enumerate}

\section{Impacto Clínico, Operacional e Normativo}
\label{sec:impacto_clinico}

A consolidação dessas ferramentas computacionais produz impacto prático imediato:
\begin{itemize}
  \item \textbf{Segurança Radiológica Personalizada:} Comprova cientificamente que reduções expressivas de dose preservam a detectabilidade diagnóstica de lesões clínicas sutis;
  \item \textbf{Auditoria e Comissionamento Hospitalar:} Fornece rotinas automatizadas para que serviços de física médica realizem auditorias em conformidade com o relatório AAPM TG-233 e com o sistema internacional IAEA 5-Star \cite{iaea_5star_2026};
  \item \textbf{Harmonização de Parques Tecnológicos:} Permite equalizar o desempenho diagnóstico entre tomógrafos de diferentes fabricantes e gerações tecnológicas.
\end{itemize}

\section{Articulação com a Pesquisa de Doutorado Direto (FAPESP 2026--2030)}
\label{sec:fapesp_doutorado}

Esta monografia de conclusão de curso cumpre o papel fundamental de consolidar o embasamento teórico, biofísico, matemático e computacional que sustenta o projeto de pesquisa de Doutorado Direto do autor (FAPESP 2026--2030) no Grupo de Dosimetria e Radioproteção em Física Médica do IFUSP:

O trabalho articula-se com os avanços em PCCT da Dra. Elsa Pimenta (Doutorado 2026) \cite{pimenta2026} e com o pipeline de automação linear no tórax de Davi Amaral (Mestrado FAPESP), estabelecendo as seguintes metas para a tese de doutorado:
\begin{enumerate}[label=\alph*)]
  \item Treinamento e validação experimental do observador DLMO baseado em \emph{Vision Transformers};
  \item Execução do estudo psicofísico nacional 2AFC com mais de 60 radiologistas especialistas sob modelagem ANOVA MRMC;
  \item Validação cruzada de transferibilidade inter-scanners (\emph{leave-one-scanner-out}) em sete tomógrafos de quatro fabricantes distintos no InRad-HCFMUSP e Radboudumc;
  \item Mapeamento experimental completo da Fronteira de Pareto Tridimensional $(D, T, -W)$ para os principais protocolos tomográficos de crânio, tórax e abdome em sistemas EICT e PCCT.
\end{enumerate}

Conclui-se, assim, este trabalho acadêmico com a convicção de que a física médica brasileira continua contribuindo para a vanguarda científica internacional, unindo o rigor analítico da física à missão de preservar vidas humanas.
